% Combinaciones marcadas en negrita y recuadro; la plantilla no admite sombreado.

\begin{table}[H]
    \centering
    \caption{Matriz morfológica, combinación A.}
    \label{tab:opcion_a}
    \footnotesize
    \setstretch{1}
    \setlength{\tabcolsep}{2.5pt}
    \renewcommand{\arraystretch}{1.15}
    \begin{tabularx}{\textwidth}{p{0.05\textwidth}>{\raggedright\arraybackslash}p{0.15\textwidth}*{3}{>{\raggedright\arraybackslash}X}}
        \toprule
        \textbf{ID} & \textbf{Función} & \textbf{Opción A} & \textbf{Opción B} & \textbf{Opción C} \\
        \midrule
        F01 & Almacenar por ingrediente & \sel{Tolvas independientes.} & Depósito compartimentado. & --- \\
        \addlinespace[2pt]
        F02 & Conducir y dosificar & \sel{Tornillos horizontales.} & Tornillos verticales descendentes. & --- \\
        \addlinespace[2pt]
        F03 & Generar el giro & \sel{Paso a paso (p.~ej., NEMA~17 o~23) y driver.} & Motorreductor CC con encoder y driver. & Servomotor industrial y accionamiento. \\
        \addlinespace[2pt]
        F04 & Transmitir el giro & \sel{Acople directo.} & Reducción por engranajes. & Correa dentada y poleas. \\
        \addlinespace[2pt]
        F05 & Distribuir el pesaje & \sel{Una balanza común.} & Balanza de harina y balanza de menores. & Receptores individuales por ingrediente. \\
        \addlinespace[2pt]
        F07 & Adquirir el peso & \sel{Convertidor dedicado (p.~ej., HX711).} & Transmisor industrial de pesaje. & Balanza con interfaz de datos. \\
        \addlinespace[2pt]
        F08 & Ejecutar la secuencia & \sel{Microcontrolador (p.~ej., ESP32).} & PLC con interfaces de pesaje. & Computador de placa única. \\
        \addlinespace[2pt]
        F09 & Ingresar consignas & \sel{Teclado y pantalla.} & Mando giratorio y pantalla. & Panel HMI. \\
        \addlinespace[2pt]
        F11 & Acondicionar la energía & \sel{Fuente AC/DC y regulación.} & Fuentes separadas por circuito. & --- \\
        \addlinespace[2pt]
        F12 & Guiar el material & \sel{Salidas orientadas al receptor.} & Colector desmontable. & --- \\
        \addlinespace[2pt]
        F13 & Vaciar el receptor auxiliar & Recipiente basculante. & Apertura inferior accionada. & \sel{Descarga directa, sin receptor auxiliar.} \\
        \addlinespace[2pt]
        F15 & Acceder para limpieza & \sel{Extracción del tornillo y del tubo.} & Retiro del módulo tolva--alimentador. & --- \\
        \bottomrule
    \end{tabularx}
    \par\smallskip
    Fuente: Elaboración propia.
\end{table}

\begin{table}[H]
    \centering
    \caption{Matriz morfológica, combinación B.}
    \label{tab:opcion_b}
    \footnotesize
    \setstretch{1}
    \setlength{\tabcolsep}{2.5pt}
    \renewcommand{\arraystretch}{1.15}
    \begin{tabularx}{\textwidth}{p{0.05\textwidth}>{\raggedright\arraybackslash}p{0.15\textwidth}*{3}{>{\raggedright\arraybackslash}X}}
        \toprule
        \textbf{ID} & \textbf{Función} & \textbf{Opción A} & \textbf{Opción B} & \textbf{Opción C} \\
        \midrule
        F01 & Almacenar por ingrediente & \sel{Tolvas independientes.} & Depósito compartimentado. & --- \\
        \addlinespace[2pt]
        F02 & Conducir y dosificar & \sel{Tornillos horizontales.} & Tornillos verticales descendentes. & --- \\
        \addlinespace[2pt]
        F03 & Generar el giro & \sel{Paso a paso (p.~ej., NEMA~17 o~23) y driver. (M)} & \sel{Motorreductor CC con encoder y driver. (H)} & Servomotor industrial y accionamiento. \\
        \addlinespace[2pt]
        F04 & Transmitir el giro & \sel{Acople directo. (M)} & \sel{Reducción por engranajes. (H)} & Correa dentada y poleas. \\
        \addlinespace[2pt]
        F05 & Distribuir el pesaje & Una balanza común. & \sel{Balanza de harina y balanza de menores.} & Receptores individuales por ingrediente. \\
        \addlinespace[2pt]
        F07 & Adquirir el peso & \sel{Convertidor dedicado (p.~ej., HX711).} & Transmisor industrial de pesaje. & Balanza con interfaz de datos. \\
        \addlinespace[2pt]
        F08 & Ejecutar la secuencia & \sel{Microcontrolador (p.~ej., ESP32).} & PLC con interfaces de pesaje. & Computador de placa única. \\
        \addlinespace[2pt]
        F09 & Ingresar consignas & \sel{Teclado y pantalla.} & Mando giratorio y pantalla. & Panel HMI. \\
        \addlinespace[2pt]
        F11 & Acondicionar la energía & \sel{Fuente AC/DC y regulación.} & Fuentes separadas por circuito. & --- \\
        \addlinespace[2pt]
        F12 & Guiar el material & \sel{Salidas orientadas al receptor. (H)} & \sel{Colector desmontable. (M)} & --- \\
        \addlinespace[2pt]
        F13 & Vaciar el receptor auxiliar & \sel{Recipiente basculante.} & Apertura inferior accionada. & Descarga directa, sin receptor auxiliar. \\
        \addlinespace[2pt]
        F15 & Acceder para limpieza & Extracción del tornillo y del tubo. & \sel{Retiro del módulo tolva--alimentador.} & --- \\
        \bottomrule
    \end{tabularx}
    \par\smallskip
    Fuente: Elaboración propia. H: harina. M: azúcar, sal y levadura.
\end{table}

\begin{table}[H]
    \centering
    \caption{Matriz morfológica, combinación C.}
    \label{tab:opcion_c}
    \footnotesize
    \setstretch{1}
    \setlength{\tabcolsep}{2.5pt}
    \renewcommand{\arraystretch}{1.15}
    \begin{tabularx}{\textwidth}{p{0.05\textwidth}>{\raggedright\arraybackslash}p{0.15\textwidth}*{3}{>{\raggedright\arraybackslash}X}}
        \toprule
        \textbf{ID} & \textbf{Función} & \textbf{Opción A} & \textbf{Opción B} & \textbf{Opción C} \\
        \midrule
        F01 & Almacenar por ingrediente & Tolvas independientes. & \sel{Depósito compartimentado.} & --- \\
        \addlinespace[2pt]
        F02 & Conducir y dosificar & Tornillos horizontales. & \sel{Tornillos verticales descendentes.} & --- \\
        \addlinespace[2pt]
        F03 & Generar el giro & \sel{Paso a paso (p.~ej., NEMA~17 o~23) y driver.} & Motorreductor CC con encoder y driver. & Servomotor industrial y accionamiento. \\
        \addlinespace[2pt]
        F04 & Transmitir el giro & \sel{Acople directo.} & Reducción por engranajes. & Correa dentada y poleas. \\
        \addlinespace[2pt]
        F05 & Distribuir el pesaje & Una balanza común. & \sel{Balanza de harina y balanza de menores.} & Receptores individuales por ingrediente. \\
        \addlinespace[2pt]
        F07 & Adquirir el peso & Convertidor dedicado (p.~ej., HX711). & \sel{Transmisor industrial de pesaje.} & Balanza con interfaz de datos. \\
        \addlinespace[2pt]
        F08 & Ejecutar la secuencia & Microcontrolador (p.~ej., ESP32). & \sel{PLC con interfaces de pesaje.} & Computador de placa única. \\
        \addlinespace[2pt]
        F09 & Ingresar consignas & Teclado y pantalla. & Mando giratorio y pantalla. & \sel{Panel HMI.} \\
        \addlinespace[2pt]
        F11 & Acondicionar la energía & Fuente AC/DC y regulación. & \sel{Fuentes separadas por circuito.} & --- \\
        \addlinespace[2pt]
        F12 & Guiar el material & \sel{Salidas orientadas al receptor. (H)} & \sel{Colector desmontable. (M)} & --- \\
        \addlinespace[2pt]
        F13 & Vaciar el receptor auxiliar & Recipiente basculante. & \sel{Apertura inferior accionada.} & Descarga directa, sin receptor auxiliar. \\
        \addlinespace[2pt]
        F15 & Acceder para limpieza & Extracción del tornillo y del tubo. & \sel{Retiro del módulo tolva--alimentador.} & --- \\
        \bottomrule
    \end{tabularx}
    \par\smallskip
    Fuente: Elaboración propia. H: harina. M: azúcar, sal y levadura.
\end{table}
